\documentclass[12pt]{article}
\usepackage[margin=1in]{geometry}
\usepackage{amsmath}
\usepackage{amssymb}
\usepackage{amsthm}
\usepackage{enumitem}
\usepackage{xcolor}
\usepackage{pagecolor}
\usepackage{marginnote}
\usepackage{fancyhdr}
\usepackage{bm}
\usepackage{etoolbox}

\definecolor{cream}{HTML}{faf7f1}
\pagecolor{cream}

\newcommand{\defn}[1]{\textbf{#1}}
\newcommand{\defeq}{\mathrel{\vcenter{\baselineskip0.5ex \lineskiplimit0pt
                     \hbox{\scriptsize.}\hbox{\scriptsize.}}}=}
\newcommand{\T}{\top}
\newcommand{\F}{\bot}
\newcommand{\Pow}{\mathbb{P}}
\newcommand{\N}{\mathbb{N}}

\newtheoremstyle{proofstyle}
  {}{}{\itshape}{}{\bfseries}{.}{ }{}
\theoremstyle{proofstyle}
\newtheorem*{proof*}{Proof}

\AtBeginEnvironment{proof*}{\setlength{\parskip}{0.5em}}

\newcommand{\sidenote}[1]{\marginnote{\small\itshape #1}}

\makeatletter
\renewcommand{\maketitle}{%
  \begin{flushleft}
    {\Large\bfseries CS173: Discrete Structures}\\[0.3em]
    {\Large\bfseries Notes}\\[0.5em]
    {\normalsize Mohammed Al Salhi}
  \end{flushleft}
  \vspace{1em}
}
\makeatother

\AtBeginDocument{\mathversion{bold}}

\begin{document}

\maketitle

\section*{Divisibility}

\textbf{Definition.} Given $x, y \in \mathbb{Z}$, we define
\[
    x \mid y \iff (\exists z \in \mathbb{Z})(x \cdot z = y).
\]
We read $x \mid y$ as ``$x$ divides $y$''. For example,
$3 \mid 6 \iff (\exists z \in \mathbb{Z})(3 \cdot z = 6) \iff 3 \cdot 2 = 6$.

\textbf{Definition.} Given $x, y \in \mathbb{Z}$, we define
\[
    x - y := x + (-y).
\]

\textbf{Definition.} We define $|\cdot| : \mathbb{Z} \to \mathbb{N}$ given by,
for every $x \in \mathbb{Z}$:
\[
    |x| :=
    \begin{cases}
        x  & \text{if } x \geq 0 \\
        -x & \text{if } x < 0
    \end{cases}
\]

\textbf{Examples.}
\begin{itemize}[itemsep=0.5em]
    \item $2 \mid 0 \iff (\exists z \in \mathbb{Z})(2 \cdot z = 0)$,
          witnessed by $z = 0$ since $2 \cdot 0 = 0$.
    \item $0 \nmid 2 \iff \neg(\exists z \in \mathbb{Z})(0 \cdot z = 2)
          \iff (\forall z \in \mathbb{Z})(0 \cdot z \neq 2)$,
          which holds since $0 \cdot z = 0 \neq 2$ for all $z$.
\end{itemize}

\section*{Divisibility is a Partial Order on $\mathbb{N}$}

\textbf{Theorem.} Divisibility on $\mathbb{N}$ is a partial order:
\begin{itemize}[itemsep=0.5em]
    \item \textbf{Reflexivity:} $(\forall x \in \mathbb{Z})(x \mid x)$.
    \item \textbf{Antisymmetry (up to sign):} $(\forall x, y \in \mathbb{Z})((x \mid y \wedge y \mid x)
          \Rightarrow |x| = |y|)$.
    \item \textbf{Transitivity:} $(\forall x, y, z \in \mathbb{Z})((x \mid y \wedge y \mid z)
          \Rightarrow x \mid z)$.
\end{itemize}

\begin{proof*}
    We prove transitivity. Let $x, y, z \in \mathbb{Z}$
    and assume $x \mid y$ and $y \mid z$. So there exist $k_1 \in \mathbb{Z}$
    and $k_2 \in \mathbb{Z}$ such that $x \cdot k_1 = y$ and $y \cdot k_2 = z$.
    Then $(x \cdot k_1) \cdot k_2 = z$, therefore $x \cdot (k_1 \cdot k_2) = z$,
    where $k_1 \cdot k_2 \in \mathbb{Z}$. Hence $x \mid z$. \qed
\end{proof*}

\textbf{Theorem.}
\begin{itemize}[itemsep=0.5em]
    \item $(\forall x \in \mathbb{Z})(x \mid 0)$. \sidenote{$0$ is a universal dividend}
    \item $(\forall x \in \mathbb{Z})(1 \mid x)$. \sidenote{$1$ is a universal divisor}
    \item $(\forall x \in \mathbb{Z})(x \mid x)$. \sidenote{reflex. of $\mid$}
    \item $(\forall x, y \in \mathbb{Z})(x \mid y \Rightarrow |x| \leq |y|)$.
\end{itemize}

\section*{Primality}

\textbf{Definition.} We say $x \in \mathbb{N}$ is \defn{prime} if
\[
    (x \geq 2) \wedge (\forall p \in \mathbb{N})(p \mid x \Rightarrow p \in \{1, x\}).
\]

\textbf{Example.}
\begin{itemize}[itemsep=0.5em]
    \item $1$ is not prime since $1 < 2$.
    \item $4$ is not prime since $2 \mid 4$ but $2 \notin \{1, 4\}$.
\end{itemize}

\textbf{Theorem.} $\mathbb{N} \subseteq \mathbb{Z}$ (false, but spiritually true).

\section*{Properties of $(\mathbb{N}, +, \cdot)$}

\begin{itemize}[itemsep=0.5em]
    \item $0$ is the $+$ identity.
    \item $+$ is associative.
    \item $+$ is commutative.
    \item $1$ is the $\cdot$ identity.
    \item $\cdot$ is associative.
    \item $\cdot$ is commutative.
    \item $\cdot$ distributes over $+$.
    \item $0$ is the annihilator of $\cdot$.
\end{itemize}

\section*{Strong Induction}

\textbf{Theorem (Strong Induction).} Let $\varphi$ be a predicate on $\mathbb{N}$.
If
\[
    \varphi(0) \wedge (\forall k \in \mathbb{N})\Big((\forall \ell \in \mathbb{N})
    (\ell < k \Rightarrow \varphi(\ell)) \Rightarrow \varphi(k)\Big)
\]
then $(\forall n \in \mathbb{N})(\varphi(n))$.

The bracketed part $(\forall \ell \in \mathbb{N})(\ell < k \Rightarrow \varphi(\ell))$
is the \defn{inductive hypothesis} --- we assume $\varphi$ holds for
\textit{all} values strictly less than $k$, not just the immediate predecessor.

\section*{Fundamental Lemma of Arithmetic}

\textbf{Theorem (Fundamental Lemma of Arithmetic).}
\[
    (\forall n \in \mathbb{N})\Big(n \geq 2 \Rightarrow (\exists p \in \mathbb{N})
    (p \text{ is prime} \wedge p \mid n)\Big)
\]

\begin{proof*}
    \textbf{Base case:} Observe $2 \mid 2$ because $2 \cdot 1 = 2$ and $1 \in \mathbb{Z}$.
    Exercise: NTS $2$ is prime.

    \textbf{Inductive step:} Let $k \in \mathbb{N}$. Assume
    \[
        (\forall \ell \in \mathbb{N})\Big(\ell < k \Rightarrow
        \big(\ell \geq 2 \Rightarrow (\exists p \in \mathbb{N})(p \text{ is prime}
        \wedge p \mid \ell)\big)\Big).
    \]
    Assume $k \geq 2$.

    \textbf{Case 1:} Suppose $k$ is prime. Then $k \mid k$, and $k$ is prime,
    so $k$ itself is the required prime divisor.

    \textbf{Case 2:} Suppose $k$ is not prime (i.e., $k$ is composite). Then
    there exists $p \in \mathbb{N}$ such that $p \mid k$ and $p \notin \{1, k\}$.
    So $p \neq 1$ and $p \neq k$. Notice $p \neq 0$ because if $p = 0$, then
    $p \mid k$ would imply $k = 0$, contradicting $k \geq 2$. So $p \geq 2$.
    Also, since $p \mid k$, we have $|p| \leq |k|$, so $p \leq k$, and since
    $p \neq k$, we get $p < k$. By the inductive hypothesis, there exists
    $q \in \mathbb{N}$ such that $q$ is prime and $q \mid p$. We know $q \mid k$
    by transitivity of divisibility.

    Therefore $(\forall n \in \mathbb{N})(n \geq 2 \Rightarrow \exists z \in \mathbb{N}
    (z \mid n \wedge z \text{ is prime}))$. \qed
\end{proof*}

\section*{Fundamental Theorem of Arithmetic}

\textbf{Theorem (Fundamental Theorem of Arithmetic).} Every $n \in \mathbb{N}$
with $n \geq 2$ has a unique prime factorization:
\[
    (\forall n \in \mathbb{N})\Big(n \geq 2 \Rightarrow
    (\exists! k \in \mathbb{N}^+)(\exists! p_1, p_2, \ldots, p_k \in \mathbb{N})
    (\exists! \alpha_1, \alpha_2, \ldots, \alpha_k \in \mathbb{N}^+)
\]
\[
    \left(n = \prod_{i=1}^{k} p_i^{\alpha_i}\right)
    \wedge (p_1 \text{ is prime}) \wedge (p_2 \text{ is prime})
    \wedge \cdots \wedge (p_k \text{ is prime})\Big)
\]

\section*{GCD and Coprimality}

\textbf{Definition (GCD).} Given $x, y \in \mathbb{Z}$, we say that $d \in \mathbb{N}$
is the \defn{greatest common divisor} of $x$ and $y$, written $\gcd(x, y) = d$, if:
\begin{enumerate}[label=\arabic*., itemsep=0.5em]
    \item $d \mid x$.
    \item $d \mid y$.
    \item $(\forall z \in \mathbb{N})(z \mid x \wedge z \mid y \Rightarrow z \mid d)$.
\end{enumerate}

\textbf{Definition (Coprime).} We say $x, y \in \mathbb{Z}$ are \defn{coprime} if
$\gcd(x, y) = 1$.

\textbf{Theorem.}
\[
    (\forall x, y \in \mathbb{Z})\Big(\gcd(x,y) = 1 \iff
    (\forall p \in \mathbb{N})\big(p \text{ is prime} \Rightarrow
    (p \nmid x) \vee (p \nmid y)\big)\Big)
\]

\textbf{Lemma.}
\begin{enumerate}[label=\arabic*., itemsep=0.5em]
    \item $(\forall x, y \in \mathbb{Z})(\gcd(x,y) = 0 \iff x = 0 \wedge y = 0)$.
    \item $(\forall x, y \in \mathbb{Z})(\gcd(x,y) = x \iff x \mid y)$.
    \item $(\forall x, y \in \mathbb{Z})(x \neq 0 \vee y \neq 0 \Rightarrow \gcd(x,y) \geq 1)$.
\end{enumerate}

\textbf{Examples.}
\begin{itemize}[itemsep=0.5em]
    \item $\gcd(1, 9) = 1$, \quad $\gcd(2, 4) = 2$, \quad $\gcd(4, 10) = 2$.
    \item $\gcd(0, 9) = 9$ since $9 \mid 9$. \quad $\gcd(4, 20) = 4 = 2^2$.
    \item $4 = 2^2$, \quad $10 = 2 \cdot 5$, \quad $20 = 2^2 \cdot 5$.
    \item $\gcd(0, 1) = 1$, \quad $(\forall x \in \mathbb{Z})(\gcd(1, x) = 1)$.
\end{itemize}

\section*{LCM}

\textbf{Definition (LCM).} The \defn{least common multiple} of $x, y \in \mathbb{Z}$
is a natural number $d \in \mathbb{N}$ such that:
\begin{enumerate}[label=\arabic*., itemsep=0.5em]
    \item $x \mid d$.
    \item $y \mid d$.
    \item $(\forall z \in \mathbb{Z})(x \mid z \wedge y \mid z \Rightarrow d \mid z)$.
\end{enumerate}

\section*{Euclid's Division Lemma}

\textbf{Theorem (Euclid's Division Lemma).} For any $x, y \in \mathbb{Z}$,
if $y \neq 0$, then
\[
    (\exists! q, r \in \mathbb{Z})(x = q \cdot y + r \ \wedge \ 0 \leq r < |y|).
\]

\textbf{Lemma.}
\[
    (\forall x, y, n \in \mathbb{Z})\Big((n \mid x \wedge n \mid y) \Rightarrow
    (\forall a, b \in \mathbb{Z})(n \mid ax + by)\Big)
\]

\textbf{Lemma.}
\[
    (\forall a, b, c \in \mathbb{Z})(a = b + c \Rightarrow \gcd(a,b) = \gcd(a,c) = \gcd(b,c))
\]

\textbf{Algorithm (Euclidean Algorithm).} For $a, b \in \mathbb{Z}$:
\[
    \gcd(a, b) :=
    \begin{cases}
        a           & \text{if } b = 0 \\
        \gcd(b, r)  & \text{if } b \neq 0,\ \text{where } r \in \mathbb{Z}
                      \text{ s.t. } (\exists q \in \mathbb{Z})(a = q \cdot b + r)
                      \wedge 0 \leq r < |b|
    \end{cases}
\]
Note: $\gcd(x, y) = \gcd(x, r)$ where $x = q \cdot y + r$ and $0 \leq r < |y|$,
so $\gcd(x,y) = \gcd(y, r)$.

\section*{Bezout's Theorem}

\textbf{Theorem (Bezout).}
\[
    (\forall x, y \in \mathbb{Z})(\exists a, b \in \mathbb{Z})(\gcd(x,y) = ax + by)
\]

\begin{proof*}
    Let $x, y \in \mathbb{Z}$. WLOG suppose $x \neq 0$ and $y \neq 0$.
    Consider
    \[
        M := \{ax + by \mid a \in \mathbb{Z} \wedge b \in \mathbb{Z} \wedge ax + by > 0\}.
    \]
    Observe $(\forall m \in M)(m \in \mathbb{Z} \wedge m > 0)$, so $M \subseteq \mathbb{N}$.
    If $x > 0$, then $x = 1 \cdot x + 0 \cdot y \in M$. If $x < 0$, then
    $-x = (-1) \cdot x + 0 \cdot y \in M$. So $M \neq \emptyset$.

    Therefore, there is some $d \in M$ such that $(\forall m \in M)(d \leq m)$.
    Since $d \in M$, we know $d = ax + by$ for some $a, b \in \mathbb{Z}$.

    We verify the three conditions for $\gcd(x, y) = d$:

    \textbf{Condition 1 --- $d \mid x$:} By Euclid's Division Lemma, there exist
    $q, r \in \mathbb{Z}$ such that $x = q \cdot d + r$ and $0 \leq r < d$. Then:
    \[
        x = q \cdot d + r \Rightarrow x = q(ax+by) + r
        \Rightarrow (1-qa)x + (-qb)y = r.
    \]
    Towards a contradiction, suppose $r \neq 0$, so $0 < r < d$. Since $r > 0$
    and $r = (1-qa)x + (-qb)y$, we know $r \in M$. But $d \leq r$ by minimality
    of $d$, contradicting $r < d$. So $r = 0$, hence $x = q \cdot d$, so $d \mid x$.

    \textbf{Condition 2 --- $d \mid y$:} By the same argument applied to $y$,
    using Euclid's Division Lemma to write $y = q' \cdot d + r'$ with $0 \leq r' < d$,
    we get $r' = (- q'a)x + (1 - q'b)y \in M$ if $r' \neq 0$, contradicting
    minimality of $d$. So $r' = 0$, hence $d \mid y$.

    \textbf{Condition 3 --- $d$ is the greatest:} Let $z \in \mathbb{N}$ and
    suppose $z \mid x$ and $z \mid y$. Since $d = ax + by$, by the linear
    combination lemma $z \mid ax + by$, so $z \mid d$.

    Therefore $d = \gcd(x, y)$, and $\gcd(x,y) = ax + by$ for some
    $a, b \in \mathbb{Z}$. \qed
\end{proof*}
\section*{Prime Divisibility}

\textbf{Theorem.}
\[
    (\forall p \in \mathbb{N})(\forall a, b \in \mathbb{Z})
    \Big((p \text{ is prime} \wedge p \mid ab) \Rightarrow (p \mid a \vee p \mid b)\Big)
\]

Note: this is special to primes. In general, $m \mid ab \not\Rightarrow (m \mid a \vee m \mid b)$.
For example, $20 \mid 50 \cdot 2$ but $20 \nmid 2$ and $20 \nmid 50$.

\end{document}